% !TeX program = xelatex
% UTF-8 transcription of six handwritten pages; source wording and doses retained.
% Source images: images/pic01.jpg through images/pic06.jpg, in page order.
% Compile directly with XeLaTeX, or input in a document using XePersian.
\ifdefined\settextfont
  \def\VZVStandalone{0}
\else
  \def\VZVStandalone{1}
  \documentclass[12pt,a4paper]{article}
  \usepackage[margin=2cm]{geometry}
  \usepackage{amsmath,amssymb,enumitem}
  \usepackage{xepersian}
  \IfFontExistsTF{Amiri}{\settextfont{Amiri}}{\settextfont{FreeSerif}}
  \setlatintextfont{FreeSerif}
\fi
\providecommand{\en}[1]{\lr{#1}}
\newcommand{\vzvpage}[2]{\clearpage\section*{\en{VZV} ـ صفحهٔ #1}%
  {\footnotesize\en{#2}\par}}
\ifnum\VZVStandalone=1 \begin{document}\fi

\vzvpage{۱}{pic01.jpg — 20261013\_034730187.jpg}
\section*{عفونت‌های \en{VZV} ـ هاریسون ۱۹۸}
\begin{itemize}
\item \en{VZV} عامل ۲ سندرم بالینی است: واریسلا (\en{chickenpox}) و هرپس زوستر (\en{shingles}).
\item معمولاً عفونت در کودکان به صورت یک بیماری خفیف با راش‌های وزیکولار اگزانتماتو بروز می‌کند.
\item سن شایع فعال شدن مجدد \en{VZV} بالای ۶۰ سال است.
\item نمای بیماری در کشت بافت:
  \begin{latin}
  Eosinophilic intranuclear inclusions and multinucleated giant cells
  \end{latin}
\end{itemize}
\subsection*{پاتوژنز و پاتولوژی}
\begin{itemize}
\item عفونت اولیه: انتقال از طریق تنفس، سپس درگیری اوروفارنکس، سپس سیستم لنفاوی، سپس ویرمی.
\item \en{PCR} خون و ضایعات مثبت است.
\item درگیری وزیکول‌ها: دژنره شدن، خونریزی، ادم؛ تجمع \en{PMN}ها محتوی وزیکول‌ها را کدر و \en{cloudy} می‌کند.
\item عفونت راجعه: مکانیسم آن ناشناخته است.
\item عفونت اولیه گانگلیون‌های ریشهٔ خلفی را درگیر می‌کند؛ گانگلیون‌های حسی دورسال، کرانیال، اتونومیک.
\item تکثیر ویروس در ارگان‌های دیگر مثل ریه ممکن است. در عفونت اولیه یا ثانویه رخ می‌دهد ولی درگیری در افراد با سیستم ایمنی ضعیف‌تر، شایع‌تر است.
\item درگیری ریه: پنومونیت بینابینی + خونریزی ریوی + انکلوزیون‌های اینترانوکلئار + تشکیل جاینت سل‌های چند هسته‌ای.
\item درگیری مغزی آن شایع است و بر خلاف \en{HSV}، \en{focal hemorrhagic necrosis} ایجاد نمی‌کند.
\end{itemize}
\subsection*{اپیدمیولوژی و تظاهرات بالینی}
\begin{itemize}
\item \en{chickenpox}: انسان تنها مخزن بیماری هست؛ نرخ انتقال ۹۰\%؛ جنس و نژاد تفاوتی ایجاد نمی‌کند.
\item ۵۰\% موارد کودکان \en{5--9} سال هستند؛ سپس زیر ۵ سال، سپس \en{10--14} سال.
\item تنها ۱۰\% موارد بالای ۱۵ سال هستند.
\item واکسیناسیون در سال دوم زندگی + قبل از مدرسه.
\end{itemize}

\vzvpage{۲}{pic02.jpg — 20261013\_034746143.jpg}
\begin{itemize}
\item طول دورهٔ کمون \en{10--21} روز ولی معمولاً \en{14--17} روز است.
\item طول دورهٔ ایجاد وزیکول \en{4--5} روز است.
\item سرایت‌پذیری از ۲ روز قبل از شروع علائم تا زمانی که وزیکول وجود دارد.
\item علائم: راش + تب \en{low-grade} + \en{malaise}.
\item گروهی از بیماران یک پرودروم \en{1--2} روزه قبل از اگزانتم دارند.
\item تب \en{37.8--39.4} برای \en{3--5} روز.
\item ضایعات پوستی: ماکولوپاپول + وزیکول + \en{scabs}.
  \begin{itemize}
  \item از تنه و صورت شروع سپس درگیری همه‌جا.
  \item هر ضایعه ظرف \en{2--4} روز نکروتیک می‌شود.
  \item می‌تواند نازوفارنکس + واژن را درگیر کند.
  \item هر چه سن کمتر، تعداد ضایعات پوستی کمتر.
  \item درگیری دوم و سوم در خانواده، تعداد ضایعات پوستی بیشتری دارند.
  \end{itemize}
\item در افراد نقص ایمنی، ضایعات مدام تا ۱۵\% نرخ مرگ و میر دارند.
% The abbreviated phrase «ضایعات مدام» is retained from the source.
\item شایع‌ترین عارضهٔ بیماری: عفونت ثانویهٔ باکتریال پوستی؛ به علت استرپتوکوک و استاف اورئوس.
\item شایع‌ترین محل درگیری خارج پوستی در کودکان، \en{CNS} است.
  \begin{itemize}
  \item سندرم آتاکسی مخچه‌ای حاد و التهاب مننژ؛ معمولاً ۲۱ روز پس از شروع راش‌ها.
  \item \en{CSF}: لنفوسیت بالا، پروتئین بالا.
  \item درگیری \en{CNS} در افراد با ایمنی سالم معمولاً خفیف است و نیاز به بستری ندارد.
  \item مننژیت آسپتیک، انسفالیت، میلیت عرضی، سندرم گیلن‌باره.
  \item شیوع انسفالیت \en{0.1--0.2} درصد بیماران آبله‌مرغان است.
  \item کودکان تحت درمان با آسپرین: سندرم \en{Reye}.
  \item هیچ درمانی (حتی آسیکلوویر) برای درگیری \en{CNS} مؤثر نیست.
  \end{itemize}
\item پنومونی واریسلا: جدی‌ترین عارضهٔ آبله‌مرغان است.
  \begin{itemize}
  \item بیشتر در بزرگسالان (تا ۲۰\% موارد).
  \item در زنان باردار بسیار شدید می‌شود.
  \item هیپوکسی، تندنفسی، سیانوز، رال‌های ریوی شایع است.
  \item در رادیولوژی درگیری انترستیشیال ندولار و پنومونیت شایع‌ترین است.
  \end{itemize}
\item سایر عوارض: میوکاردیت، ضایعات \en{Cornea}، نفریت، آرتریت، \en{Bleeding Diathesis}، گلومرولونفریت حاد، هپاتیت.
\item هپاتیت معمولاً بدون علامت است و افزایش آنزیم وجود دارد: \en{AST, ALT}.
\end{itemize}

\vzvpage{۳}{pic03.jpg — 20261013\_034754002.jpg}
\subsection*{واریسلای پری‌ناتال}
\begin{itemize}
\item تا ۳۰\% نرخ مرگ و میر.
\item چنانچه مادر ظرف ۵ روز قبل تا ۲ روز بعد از زایمان به بیماری مبتلا شود، بیماری به شدت شدید بروز می‌کند.
\end{itemize}
\subsection*{هرپس زوستر}
\begin{itemize}
\item در هر سنی ممکن است دیده شود؛ بیشتر از دههٔ ۶۰ به بعد.
\item واکسیناسیون نرخ بروز \en{shingles} را کاهش داده است.
\item زونای راجعه بسیار بسیار نادر است مگر \en{AIDS}.
\item شایع‌ترین درماتوم‌های درگیر \en{$T_3$--$L_3$} می‌باشند.
\item درگیری شاخهٔ افتالمیک: زوستر افتالمیکوس.
\item زونا در کودکان خفیف است ولی در بزرگسالان شدیدتر است.
\item درد \en{herald}: \en{48--72} ساعت قبل از ضایعات پوستی.
\item راش ماکولوپاپولر اریتماتو، سپس وزیکول.
\item کل طول دورهٔ بیماری معمولاً \en{7--10} روز است؛ گاه تا \en{2--4} هفته.
\item می‌تواند ویروس را به افراد غیر ایمن انتقال دهد: \en{chickenpox}.
\item درد در درماتوم + ویرمی بدون ضایعات پوستی: \en{Zoster sine herpetica}.
% Latin spelling is preserved as written in the notes.
\item درگیری تری‌ژمینال: ضایعات در نوک بینی، \en{Hutchinson Sign}؛ سمت دهان و حفرهٔ نازال.
\item درگیری افتالمیک: در غیاب درمان ضد ویروسی، کوری.
\item سندرم \en{Ramsay-Hunt}:
  \begin{itemize}
  \item درد و وزیکول در کانال گوش خارجی.
  \item از دست دادن چشایی در \en{$\frac{2}{3}$} قدامی زبان.
  \item فلج صورت همان سمت.
  \item درگیری گانگلیون ژنیکولت شاخهٔ حسی عصب \en{facial}.
  \end{itemize}
\item پیامد هرپس نورالژیا: در بیماران جوان نادر است؛ ۵۰\% بیماران بالای ۵۰ سال تا حدودی درد تا ماه‌ها دارند؛ هایپو و هایپراستزی شایع است.
% Source reads «پیامد هرپس نورالژیا».
\item ممکن است درگیری \en{CNS} بدون علامت بروز کند: \en{CSF+}.
\item مننگوانسفالیت: سردرد، تب، فوتوفوبیا، مننژیت، استفراغ.
\item بیماران هوچکین و لنفوم \en{non-Hodgkin}: بیماری شدید.
\end{itemize}

\vzvpage{۴}{pic04.jpg — 20261013\_034803300.jpg}
\subsection*{تشخیص‌های افتراقی}
\begin{itemize}
\item تشخیص آبله‌مرغان بر اساس تظاهرات بالینی است: وجود راش کلاسیک + سابقهٔ مواجهه با فرد مبتلا.
\item \en{DDX}:
  \begin{itemize}
  \item عفونت \en{HSV} منتشر در بیمار مبتلا به درماتیت آتوپیک.
  \item برخی عفونت‌های کوکساکی ویروس.
  \item سرخک (\en{measles}) آتیپیک.
  \item برخی عفونت‌های اکوویروس.
  \end{itemize}
\item افتراق از \en{Smallpox}: ضایعات در آبله یکنواخت هستند و همگی در یک مرحله قرار دارند.
\item \en{HSV} و کوکساکی ویروس‌ها می‌توانند ضایعات وزیکولار درماتومال ایجاد کنند.
\end{itemize}
\subsection*{یافته‌های آزمایشگاهی}
\begin{itemize}
\item تشخیص قطعی صرفاً با ایزوله کردن ویروس قابل انجام است: کشت، \en{PCR}.
\item ۴ برابر شدن سطح \en{Ab} هم تشخیص قطعی است.
\item برای \en{Rapid Test}: خون، ضایعات، بزاق.
\item رایج‌ترین تست سرولوژیک: \en{ELISA, FAMA}.
\end{itemize}
\subsection*{درمان ویروس واریسلا زوستر}
\subsubsection*{آبله‌مرغان}
\begin{itemize}
\item شستن روزانهٔ بدن، شستن مرتب دست‌ها، محدودیت تماس در قرنطینه، تجویز داروی ضد خارش.
\item تجویز آسپرین برای کودکان مبتلا به آبله‌مرغان ممنوع است: \en{Reye Synd}.
\item ظرف ۲۴ ساعت اولیهٔ بیماری در نوجوانان و بزرگسالان، اندیکاسیون دارو:
  \begin{itemize}
  \item آسیکلوویر \en{800}، ۵ بار در روز.
  \item والاسیکلوویر \en{1 gr}، ۳ بار در روز.
  \item فامسیکلوویر \en{250}، ۳ بار در روز.
  \item مدت: \en{5--7} روز.
  \end{itemize}
\item آسیکلوویر اطفال در سن زیر ۱۲ سال هم مؤثر است: \en{20 mg/kg q6h}.
\end{itemize}
\subsubsection*{زونا}
\begin{itemize}
\item مهارکننده‌های آنزیم [خوانش نامطمئن: انزیمات].
% The enzyme phrase is unclear in the handwritten source.
\item درمان ضد ویروس خوراکی می‌دهیم.
\item آسیکلوویر \en{800}، ۵ بار، \en{7--10} روز.
\item والا و فامسیکلوویر بهتر هستند و ارجح می‌باشند:
  \begin{itemize}
  \item والاسیکلوویر \en{1 gr TDS}، \en{7--10} روز.
  \item فامسیکلوویر \en{500 TDS}.
  \end{itemize}
\end{itemize}

\vzvpage{۵}{pic05.jpg — 20261013\_034807881.jpg}
\begin{itemize}
\item در بیماران با نقص ایمنی شدید، آبله‌مرغان و زونا باید با آسیکلوویر \en{IV} درمان شوند: \en{10 mg/kg q8h for 7 days}.
\item بیماران با پنومونی واریسلا معمولاً نیاز به \en{MV} پیدا می‌کنند.
\item درگیری چشمی زونا: ارجاع به چشم‌پزشک، آتروپین.
\end{itemize}
\subsection*{درمان نوریت و نورالژی پس از هرپس}
\begin{itemize}
\item گاباپنتین، پرگابالین، آمی‌تریپتیلین، لیدوکائین (پچ)، فلوفنازین.
\item پردنیزون:
  \begin{itemize}
  \item \en{60 mg} روزانه در روز \en{1--7}.
  \item \en{30 mg} روزانه در روز \en{8--14}.
  \item \en{15 mg} روزانه در روز \en{15--21}.
  \end{itemize}
\item در استئوپروز، \en{DM}، گلوکوم، \en{HTN} نمی‌دهیم.
\item بدون ضد ویروس نمی‌دهیم.
\end{itemize}
\subsection*{پیشگیری}
\begin{itemize}
\item واکسن زندهٔ واریسلا (\en{Oka}): در کودکان بالای ۱ سال (تا ۱۲ سال) که تاکنون مبتلا نشده‌اند.
  \begin{itemize}
  \item اگر سن بین \en{1--2} سال: دوز دوم قبل از مدرسه، \en{12--15} ماهگی / \en{4--6} سالگی.
  \item برای همهٔ بزرگسالان سرونگاتیو هم توصیه می‌شود: ۲ دوز با فاصلهٔ حداقل ۱ ماه.
  \end{itemize}
\item واکسن شینگریکس: برای بزرگسالان بالای ۵۰ سال توصیه می‌شود.
  \begin{itemize}
  \item دو دوز با فاصلهٔ حداقل ۱ ماه.
  \item تا \en{97.2} درصد از زونا جلوگیری می‌کند.
  \item ایمنی واکسن در طول زمان کاهش می‌یابد.
  \end{itemize}
\item ایمونوگلوبولین واریسلا زوستر (\en{VZIG}):
  \begin{itemize}
  \item برای افرادی که مستعد ابتلا هستند + ریسک بالای ایجاد عارضه دارند + مواجههٔ جدی داشته‌اند.
  \item حداکثر ظرف \en{96 h} (بهتر \en{72 h}) باید تزریق شود.
  \item ولی تا ۱۰ روز هم می‌توان تزریق کرد.
  \end{itemize}
\item داروی ضد ویروس پروفیلاکسی: افرادی که \en{high-risk} هستند ولی اندیکاسیون واکسن ندارند و از ۹۶ ساعت تزریق \en{VZIG} هم گذشته است.
\item شروع درمان از روز ۷ام بعد از مواجهه.
\end{itemize}

\vzvpage{۶}{pic06.jpg — 20261013\_034816352.jpg}
\subsection*{توصیه‌های بالینی تزریق \en{VZIG}}
\subsubsection*{کرایتریای مواجهه}
\begin{enumerate}
\item مواجههٔ جدی با بیمار مبتلا به آبله‌مرغان یا زونا:
  \begin{itemize}
  \item افراد مقیم یک خانه.
  \item هم‌بازی‌های کودک در یک محیط سرپوشیده.
  \item در بیمارستان: \en{2--4} تخت کناری در آبله‌مرغان؛ تماس مستقیم با فرد بیمار در زونا.
  \item نوزاد: مادر از ۵ روز قبل تا ۲ روز بعد از زایمان مبتلا شده باشد، آبله‌مرغان.
  \item برای زونا اندیکاسیون ندارد.
  \end{itemize}
\item \en{VZIG} باید در اولین فرصت داده شود ولی نه بعد از ۹۶ ساعت.
% The 96-hour statement here and the 10-day statement on page 5 both occur in the source.
\end{enumerate}
\subsubsection*{کاندیدها}
\begin{enumerate}
\item کودکان نقص ایمنی مستعد ابتلا که سابقهٔ واریسلا یا تزریق واکسن ندارند.
\item زنان باردار مستعد.
\item نوزادانی که مادری با شروع آبله‌مرغانش ظرف ۵ روز قبل تا ۲ روز بعد از زایمان بوده است.
\item نوزادان نارس بستری (\en{$GA\geq28$} هفته) که مادرشان سابقهٔ قابل اعتمادی از ابتلا یا واکسن ندارد و یا سرولوژی منفی دارد.
\item نوزادان نارس بستری (\en{$GA\leq28$} هفته یا \en{$BW\leq1000$ gr}) فارغ از سابقه یا سرولوژی مادر.
% Inequality signs above follow the handwritten source.
\end{enumerate}

\ifnum\VZVStandalone=1
  \let\VZVEnd\enddocument
\else
  \let\VZVEnd\relax
\fi
\VZVEnd
